% 02_related_work.tex -- DRAFT (TV3). Assumes natbib (\citep / \citet) and an English ACL-style template.
% Every citation below was checked against its paper or ACL Anthology page; see related_work_refs.bib
% for the entries that still carry a TODO. Replace [TODO: ...] markers before submission.

\section{Related Work}
\label{sec:related}

\paragraph{Vietnamese NER resources.}
Vietnamese NER has long relied on the VLSP 2016 and VLSP 2018 shared tasks
\citep{nguyen2016vlsp,nguyen2018vlsp}, which cover generic entity types
(persons, organizations, locations) in news text.
\citet{truong2021covid} introduced PhoNER\_COVID19, a domain-specific dataset for
COVID-19 patient reports with 10 entity types and about 35K entities over 10K sentences.
The annotation was done on syllable-level text, and a word-level variant was
obtained with RDRSegmenter \citep{nguyen2018rdr}, which is part of VnCoreNLP \citep{vu2018vncorenlp}.
The benchmark systems are BiLSTM-CNN-CRF \citep{ma2016endtoend} and the pre-trained
encoders XLM-R \citep{conneau2020xlmr} and PhoBERT \citep{nguyen2020phobert};
the best system, PhoBERT$_{\text{large}}$, reaches 0.945 micro-F1 on the word-level data, and word
segmentation was found to help.
Their error analysis shows that the hardest cases are the \textsc{location}/\textsc{organization}
ambiguity and over-generation of generic place mentions.
Later work treats NER as text-to-text generation; for instance, ViT5 reports
94.5 micro-F1 on PhoNER\_COVID19 \citep{phan2022vit5}.
These supervised systems are our reference point (Section~\ref{sec:experiments}).

\paragraph{NER with large language models.}
Sequence labeling does not fit the output format of a generative LLM directly.
GPT-NER \citep{wang2023gptner} therefore casts NER as generation: the model
copies the input and marks entity positions, and the output is mapped back to labels.
On five widely used benchmarks it reaches results comparable to supervised baselines and
is particularly strong when labeled data is scarce.
\citet{xie2023empirical} present an empirical study of zero-shot NER with ChatGPT.
Our setting differs in two ways: the text is Vietnamese and domain-specific, and we
ask the model for a list of entity strings with types instead of a marked copy of the sentence.
This list is aligned to the tokens by a deterministic procedure
(Section~\ref{sec:method-llm}), so every failure of the model (invented strings,
strings that cross token boundaries, ill-formed JSON) can be counted.

\paragraph{LLMs for Vietnamese and medical NER.}
Evidence for Vietnamese is still thin.
In the construction of the VietMed-NER spoken medical corpus, \citet{vietmed2024} report that
prompting GPT-4 to pre-tag entities did not reach acceptable accuracy for their entity types.
[TODO: add 1--2 recent Vietnamese LLM-for-NER or Vietnamese LLM benchmark papers after a
targeted literature search; none were found in the search done for this draft.]

\paragraph{Position of this work.}
We compare a prompted LLM (Gemini, zero-shot [TODO: and few-shot]) with the supervised
baselines of \citet{truong2021covid} on PhoNER\_COVID19 under the same strict
entity-level scoring, and we report the rate of malformed outputs and unmatched entities
alongside F1 [TODO: insert numbers from the Week 4 runs].
The annotation guidelines contain context-dependent rules, for example that a hospital is a
\textsc{location} when the patient was there and an \textsc{organization} when it acts as an agent,
and that occupations, ages and genders are tagged only when linked to a named person or patient ID;
we examine whether a zero-shot model can follow such rules from the prompt alone.
